\chapter{Polyphase propagation and scheduling}\label{ch:phases}
Storage blocks say where data are held. Structural phases say where an
operator can send a vector. Exact zero couplings produce finite propagation
frontiers. Weak couplings produce a controlled leakage problem instead.
The scheduler must distinguish those two cases and state what knowledge of
the phase graph is assumed.
\section{Polyphase structure}\label{sec:final-14}

Let \(P_h\) project onto variant phase \(h\), and again set
\[
H=A^\top A.
\]
Define the exact phase graph by
\[
g\to h
\quad\Longleftrightarrow\quad
P_hHP_g\ne0.
\]

\begin{theorem}[Polyphase light cone]\label{claim:final-14-1}
If the graph distance from phase \(g\) to phase \(h\) exceeds \(k\), then
\[
\boxed{
P_hH^kP_g=0.
}
\]
\end{theorem}

\begin{proof}
Insert the projector resolution
\(\sum_jP_j=I\) between successive factors of \(H\).  Every nonzero term in
\(P_hH^kP_g\) corresponds to a phase path of length \(k\).  If no such path
exists, every term vanishes.
\end{proof}


\begin{figure}[htbp]
\centering
\includegraphics[width=.82\linewidth]{figure_04_polyphase_lightcone.png}
\caption{The exact polyphase light cone. The support of the Krylov iterate expands by at most one graph edge per operator application. Phases outside graph distance $k$ remain identically inactive at depth $k$ and need not be loaded.}
\label{fig:polyphase-lightcone}
\end{figure}

\begin{theorem}[Streaming-polyphase composition]\label{claim:final-14-2}
At Krylov depth \(k\), an exact matrix-free iteration may omit every phase
outside the exact light cone.  Within each active phase, arbitrary storage
refinement may be used.
\end{theorem}

\begin{proof}
The light-cone theorem removes inactive phase contributions exactly.
Partition-refinement invariance shows that the remaining active phase
operators may be streamed through any exact subpartition.
\end{proof}

This is the key composition law:
\[
\boxed{
\text{mathematical locality}
\longrightarrow
\text{active phases}
\longrightarrow
\text{active storage blocks}.
}
\]


\section{Certified approximate light cones}\label{sec:final-15}
\label{sec:approx-lightcone}

The exact light-cone theorem requires literal zero phase couplings.  Real LD,
annotation, and feature graphs are often only approximately sparse.  This
section replaces exact support by a computable phasewise envelope.

Let the phase projectors be $P_1,\ldots,P_q$, and let $X$ be a vector or block
vector.  Define its phase-norm profile
\[
\rho(X)_g=\lVert P_gX\rVert_F,
\]
and define the nonnegative phase-transfer matrix
\[
\mathcal B_{hg}=\lVert P_hHP_g\rVert_2.
\]
The matrix $\mathcal B$ records the strength of every one-step phase transfer
without assuming that weak couplings vanish.

\begin{theorem}[Weighted phase-transfer envelope]
\label{thm:weighted-transfer}
For every block vector $X$,
\[
\boxed{
\rho(HX)\leq \mathcal B\rho(X)
}
\]
componentwise.  Consequently, for every integer $k\geq0$,
\[
\boxed{
\rho(H^kX)\leq \mathcal B^k\rho(X).
}
\]
\end{theorem}

\begin{proof}
Using $\sum_gP_g=I$ and submultiplicativity,
\[
\begin{aligned}
\lVert P_hHX\rVert_F
&=\left\lVert\sum_gP_hHP_gX\right\rVert_F\\
&\leq\sum_g\lVert P_hHP_g\rVert_2\lVert P_gX\rVert_F\\
&=(\mathcal B\rho(X))_h.
\end{aligned}
\]
Iteration gives the power bound.
\end{proof}

Now choose a symmetric set of retained phase blocks and define
\[
H_0=\sum_{(h,g)\in\mathcal E_0}P_hHP_g.
\]
Let
\[
(\mathcal B_0)_{hg}=
\begin{cases}
\mathcal B_{hg},&(h,g)\in\mathcal E_0,\\
0,&\text{otherwise},
\end{cases}
\qquad
\mathcal R=\mathcal B-\mathcal B_0.
\]
Symmetric pair retention preserves self-adjointness of $H_0$, although the
following error theorem does not require positive semidefiniteness of $H_0$.

\begin{theorem}[Thresholded-power certificate]
\label{thm:thresholded-power}
For every $k\geq1$,
\[
\boxed{
\rho\bigl((H^k-H_0^k)X\bigr)
\leq
(\mathcal B^k-\mathcal B_0^k)\rho(X).
}
\]
In particular,
\[
\boxed{
\lVert(H^k-H_0^k)X\rVert_F
\leq
\left\lVert(\mathcal B^k-\mathcal B_0^k)\rho(X)\right\rVert_2.
}
\]
\end{theorem}

\begin{proof}
Write $E=H-H_0$.  The telescoping identity gives
\[
H^k-H_0^k
=
\sum_{j=0}^{k-1}H^{k-1-j}EH_0^j.
\]
The block-norm matrix of $E$ is $\mathcal R$.  Applying
Theorem~\ref{thm:weighted-transfer} successively to each term gives
\[
\rho\bigl((H^k-H_0^k)X\bigr)
\leq
\sum_{j=0}^{k-1}
\mathcal B^{k-1-j}\mathcal R\mathcal B_0^j\rho(X).
\]
The matrix telescoping identity
\[
\mathcal B^k-\mathcal B_0^k
=
\sum_{j=0}^{k-1}
\mathcal B^{k-1-j}(\mathcal B-\mathcal B_0)\mathcal B_0^j
\]
yields the phasewise formula.  Orthogonality of the phase projectors gives
\[
\lVert Y\rVert_F^2=\sum_h\lVert P_hY\rVert_F^2,
\]
which proves the global bound.
\end{proof}

\begin{corollary}[Certified approximate light cone]
\label{cor:approx-lightcone}
Assume $X$ is supported on a set of start phases $S$.  If phase $h$ is not
reachable from $S$ by a retained path of length at most $k$, then
\[
P_hH_0^kX=0
\]
and
\[
\boxed{
\lVert P_hH^kX\rVert_F
\leq
\bigl[(\mathcal B^k-\mathcal B_0^k)\rho(X)\bigr]_h.
}
\]
Thus an exact zero outside the retained light cone is replaced by a certified
nonzero leakage budget.
\end{corollary}

\begin{corollary}[Scalar threshold law]\label{claim:final-15-4}
For any induced matrix norm, let
\[
M=\max\{\lVert\mathcal B\rVert,\lVert\mathcal B_0\rVert\}.
\]
Then
\[
\boxed{
\lVert\mathcal B^k-\mathcal B_0^k\rVert
\leq
kM^{k-1}\lVert\mathcal R\rVert.
}
\]
Hence weak discarded coupling produces a first-order error bound proportional
to the discarded transfer norm and to the Krylov depth.
\end{corollary}

\begin{theorem}[Krylov-polynomial certificate]
\label{thm:polynomial-certificate}
Let
\[
p(t)=\sum_{j=0}^d a_jt^j.
\]
Then
\[
\boxed{
\rho\bigl((p(H)-p(H_0))X\bigr)
\leq
\sum_{j=1}^d
|a_j|(\mathcal B^j-\mathcal B_0^j)\rho(X).
}
\]
\end{theorem}

\begin{proof}
Apply Theorem~\ref{thm:thresholded-power} to every nonconstant power and use
the triangle inequality.  The constant terms cancel.
\end{proof}

The polynomial theorem covers arbitrary vectors in a finite Krylov space,
including polynomial filters and restarted combinations, rather than only the
monomial $H^kX$.

\begin{corollary}[Certified phase pruning]\label{claim:final-15-6}
If a thresholded phase graph is chosen so that
\[
\left\lVert
(\mathcal B^k-\mathcal B_0^k)\rho(X)
\right\rVert_2
\leq\tau,
\]
then replacing $H^kX$ by $H_0^kX$ incurs Frobenius error at most $\tau$.
The same statement holds for a Krylov polynomial with the certificate in
Theorem~\ref{thm:polynomial-certificate}.
\end{corollary}

This turns graph thresholding into an error-budget decision.  A phase or
interface is not discarded because it merely appears small. It is discarded
only when the weighted sum of all paths using discarded couplings remains
inside the requested tolerance.

\begin{figure}[htbp]
\centering
\includegraphics[width=.84\linewidth]{figure_09_approximate_lightcone_certificate.png}
\caption{Weak long-range phase couplings replace the exact zero outside the retained light cone by a nonzero but certified envelope.  In this PSD test, the phasewise transfer-matrix certificate bounds the observed leakage at every depth from one through eight.}
\label{fig:approx-lightcone-certificate}
\end{figure}

\begin{figure}[htbp]
\centering
\includegraphics[width=.84\linewidth]{figure_10_threshold_tradeoff.png}
\caption{Certified phase-graph sparsification at Krylov depth five.  Retaining fewer off-diagonal phase blocks reduces communication and operator work while the transfer-matrix bound remains above the measured truncation error.  The bound is intentionally conservative because it ignores cancellation and singular-vector alignment.}
\label{fig:threshold-tradeoff}
\end{figure}

\subsection{Numerical audit}

The validation operator has twelve phases of width four.  A block-bidiagonal
factor generates a positive-definite local operator $H_0$, and weak symmetric
long-range blocks are then added to obtain $H$.  Both normalized operators
remain positive definite.  A three-column start block is supported in phase
zero.

For depths one through eight:
\begin{itemize}
\item every observed phase norm lay below $\mathcal B^k\rho(X)$.
\item every observed thresholding error lay below
$(\mathcal B^k-\mathcal B_0^k)\rho(X)$;
\item the largest total observed-error-to-certificate ratio was approximately
$0.669$.
\end{itemize}

A separate hard-threshold sweep retained between $22$ and $132$ directed
off-diagonal phase blocks.  Every measured depth-five error remained below its
certificate. The largest observed-error-to-certificate ratio was approximately
$0.376$.  Finally, a degree-four signed polynomial in $H$ satisfied the
phasewise polynomial certificate in all twelve phases, with maximum ratio
approximately $0.068$.

These computations validate the implementation and illustrate sharpness.  The
universal result is supplied by the preceding proofs, not by the finite test.



\section{Certified budget-optimal phase scheduling}\label{sec:final-16}
\label{sec:budget-scheduler}

The preceding section answers a validation question: given a retained phase
operator $H_0$, how large can the Krylov error be?  This section turns that
certificate into a selection rule.  The goal is to remove as much phase
communication as possible while remaining inside a prescribed error budget.

Let $\mathcal E$ denote the set of undirected off-diagonal phase pairs.  Because
$H$ is self-adjoint, $\mathcal B$ is symmetric.  For a phase pair
$e=\{h,g\}$ define the nonnegative matrix $\Delta_e$ by
\[
(\Delta_e)_{hg}=\mathcal B_{hg},\qquad
(\Delta_e)_{gh}=\mathcal B_{gh},
\]
with every other entry zero.  If $\mathcal O\subseteq\mathcal E$ is the set
of omitted pairs, define
\[
\mathcal B_{\mathcal O}
=\mathcal B-\sum_{e\in\mathcal O}\Delta_e.
\]
The diagonal is always retained.

For depth $k\ge1$ and starting phase profile $\rho_0=\rho(X)$, define the
\emph{edge risk vector}
\[
\boxed{
q_e^{(k)}
=
\sum_{j=0}^{k-1}
\mathcal B^{k-1-j}\Delta_e\mathcal B^j\rho_0
}
\]
and its scalar risk
\[
\boxed{
\beta_e^{(k)}=\lVert q_e^{(k)}\rVert_2.
}
\]
Unlike the raw block norm $\mathcal B_{hg}$, this quantity depends on the
start profile and on Krylov depth.

\begin{theorem}[Edge-additive leakage majorant]
\label{thm:edge-additive-majorant}
For every omitted set $\mathcal O\subseteq\mathcal E$,
\[
\boxed{
0\le
(\mathcal B^k-\mathcal B_{\mathcal O}^k)\rho_0
\le
\sum_{e\in\mathcal O}q_e^{(k)}
}
\]
componentwise.  Consequently,
\[
\boxed{
\lVert(H^k-H_{\mathcal O}^k)X\rVert_F
\le
\left\lVert
(\mathcal B^k-\mathcal B_{\mathcal O}^k)\rho_0
\right\rVert_2
\le
\sum_{e\in\mathcal O}\beta_e^{(k)}.
}
\]
\end{theorem}

\begin{proof}
The first inequality in the second display is
Theorem~\ref{thm:thresholded-power}.  For the transfer matrices, the exact
noncommutative telescoping identity is
\[
\mathcal B^k-\mathcal B_{\mathcal O}^k
=
\sum_{j=0}^{k-1}
\mathcal B^{k-1-j}
(\mathcal B-\mathcal B_{\mathcal O})
\mathcal B_{\mathcal O}^j.
\]
All matrices and $\rho_0$ are nonnegative, and
$0\le\mathcal B_{\mathcal O}\le\mathcal B$ entrywise.  Hence
\[
\mathcal B_{\mathcal O}^j\rho_0
\le
\mathcal B^j\rho_0.
\]
Substituting
\[
\mathcal B-\mathcal B_{\mathcal O}
=
\sum_{e\in\mathcal O}\Delta_e
\]
and distributing the sum gives the componentwise majorant.  The final scalar
bound follows from the Euclidean triangle inequality.
\end{proof}

\begin{corollary}[Polynomial edge-risk certificate]\label{claim:final-16-2}
Let
\[
p(t)=\sum_{d=0}^{D}a_dt^d.
\]
Define
\[
q_e^{(p)}
=
\sum_{d=1}^{D}|a_d|q_e^{(d)},
\qquad
\beta_e^{(p)}=\lVert q_e^{(p)}\rVert_2.
\]
Then
\[
\boxed{
\lVert(p(H)-p(H_{\mathcal O}))X\rVert_F
\le
\sum_{e\in\mathcal O}\beta_e^{(p)}.
}
\]
Thus the same scheduling machinery applies to finite Krylov combinations and
polynomial filters.
\end{corollary}

\begin{corollary}[Query-invisible phase pair]\label{claim:final-16-3}
If
\[
\Delta_e\mathcal B^j\rho_0=0
\qquad
\text{for every }0\le j<k,
\]
then
\[
\beta_e^{(k)}=0.
\]
Therefore a phase pair can be large in operator norm and still have zero
risk for a particular finite-depth query.  A magnitude-only threshold cannot
represent this distinction.
\end{corollary}

\subsection{Equal-cost optimality}

Suppose each retained phase pair has the same communication or scheduling
cost.  Under the additive certificate, a set of omitted pairs is feasible if
\[
\sum_{e\in\mathcal O}\beta_e^{(k)}\le\tau.
\]
Sort the risks as
\[
\beta_{(1)}\le\beta_{(2)}\le\cdots\le\beta_{(N)}.
\]

\begin{theorem}[Cardinality-optimal certified pruning]
\label{thm:unit-cost-optimal}
Let $r$ be the largest integer satisfying
\[
\sum_{i=1}^{r}\beta_{(i)}\le\tau.
\]
Then omitting the $r$ least-risk phase pairs is feasible, and no feasible
omitted set can contain more than $r$ pairs under the scalar additive
certificate.
\end{theorem}

\begin{proof}
Among all subsets of cardinality $r+1$, the smallest possible sum of risks is
\[
\sum_{i=1}^{r+1}\beta_{(i)}.
\]
By maximality of $r$, that sum exceeds $\tau$.  Hence no set of $r+1$ omitted
pairs is feasible.  The prefix of length $r$ is feasible by construction.
\end{proof}

This is a genuine optimization theorem, not a heuristic threshold rule.  The
optimization target is the additive certificate itself. The tighter exact
transfer envelope can leave additional slack.

\subsection{Exact-envelope refinement}

After the cardinality-optimal additive schedule is found, recomputing
\[
\left\lVert
(\mathcal B^k-\mathcal B_{\mathcal O}^k)\rho_0
\right\rVert_2
\]
may reveal unused budget.  Additional pairs can then be removed one at a time,
provided the exact transfer envelope remains below $\tau$.

\begin{proposition}[Certified refinement]\label{claim:final-16-5}
Any sequence of additional omissions accepted only when
\[
\left\lVert
(\mathcal B^k-\mathcal B_{\mathcal O}^k)\rho_0
\right\rVert_2\le\tau
\]
remains a valid $\tau$-certificate for the original operator error.
\end{proposition}

The refinement rule used in the reference implementation chooses, among all
currently feasible next omissions, the pair that produces the smallest
increase in the exact transfer envelope.  It is certified but is not claimed
globally optimal.

\subsection{Unequal communication costs}

Let $c_e\ge0$ denote the communication or storage cost saved by omitting pair
$e$.  Under the scalar additive certificate, the exact discrete selection
problem is
\[
\boxed{
\max_{x_e\in\{0,1\}}
\sum_e c_ex_e
\quad\text{subject to}\quad
\sum_e\beta_e^{(k)}x_e\le\tau.
}
\]
This is the $0$--$1$ knapsack problem.  Hence unequal-cost scheduling is
combinatorial in general.  The implementation supplies a safe
cost-to-risk greedy heuristic. Its output remains certified because it never
exceeds the additive risk budget.  A fractional-knapsack relaxation supplies
an upper bound on the maximum possible saved cost.

\begin{figure}[htbp]
\centering
\includegraphics[width=.84\linewidth]{figure_11_budget_scheduler_comparison.png}
\caption{Depth-five pruning under a fixed certified error budget.  A global block-norm threshold is nearly unable to prune this example without crossing the requested tolerance, whereas query-aware edge risks remove many phase pairs.  The exact-envelope refinement safely uses slack left by the additive majorant.}
\label{fig:budget-scheduler-comparison}
\end{figure}

\begin{figure}[htbp]
\centering
\includegraphics[width=.82\linewidth]{figure_13_edge_risk_vs_block_norm.png}
\caption{Raw phase-block magnitude and finite-depth query risk are different quantities.  The edge score $\beta_{hg}^{(5)}$ incorporates the start profile and every depth-five transfer path through the phase graph. The block norm $\mathcal B_{hg}$ does not.}
\label{fig:edge-risk-vs-norm}
\end{figure}

\subsection{Numerical audit}

The twelve-phase operator from Section~\ref{sec:approx-lightcone} was reused at
Krylov depth five.  There are $66$ undirected off-diagonal phase pairs.
For error budgets $10^{-3}$, $2\times10^{-3}$, $5\times10^{-3}$, and
$10^{-2}$, the additive-optimal scheduler retained respectively
$61$, $56$, $45$, and $36$ pairs.  Exact-envelope refinement reduced those
counts further to
$58$, $50$, $36$, and $30$ pairs while remaining inside budget.

In the same test, the best single global threshold on the raw block norm
retained $65$ pairs at every one of those four budgets.  The difference is
structural: the global threshold is blind to start phase and Krylov depth.

A separate six-phase problem contains $15$ candidate phase pairs, so all
$2^{15}=32768$ subsets can be enumerated.  At four distinct budgets, the
sorted-risk scheduler matched the brute-force cardinality optimum under the
additive certificate in every case.  The exact-envelope greedy refinement also
matched the brute-force exact-envelope optimum in cardinality for all four
budgets in this finite audit.  The latter agreement is empirical and is not
promoted to a theorem.

For an unequal-cost eight-phase test, costs were proportional to the bytes in
symmetric dense phase blocks.  The cost-to-risk heuristic saved $5408$ bytes
out of a $9472$-byte communication proxy while using additive risk $0.011632$
under budget $0.012$.  The fractional relaxation upper bound was $5647.8$
bytes, placing the greedy schedule within about $4.25\%$ of that relaxation
bound on this instance.


\section{Asynchronous double-buffer execution model}\label{sec:final-17}
\label{sec:double-buffer}

The mathematical streaming identities do not require transfers and GEMMs to be
serial.  Consider $B$ storage blocks.  Let $t_b$ be the transfer time for block
$b$ and $c_b$ its device compute time.

\begin{proposition}[Ideal two-buffer pipeline]\label{claim:final-17-1}
In a model with perfect asynchronous overlap, transfer of block $b+1$ may run
while block $b$ is being processed.  The serial one-buffer time is
\[
T_{\rm serial}=\sum_{b=1}^{B}(t_b+c_b),
\]
whereas the two-buffer pipeline time is
\[
\boxed{
T_{\rm pipe}
=
t_1
+
\sum_{b=1}^{B-1}\max\{c_b,t_{b+1}\}
+
c_B.
}
\]
For uniform times $t_b=t$ and $c_b=c$,
\[
T_{\rm pipe}=t+c+(B-1)\max\{t,c\}.
\]
\end{proposition}

\begin{proof}
After the first block arrives, every interior pipeline stage finishes only when
both the current computation and the next transfer have completed.  Their
joint stage duration is therefore the maximum of those two times.  The initial
transfer and final compute cannot be hidden and form the two boundary terms.
\end{proof}

If one resident block uses $M_{\rm block}$ bytes and all other fixed working
state uses $M_{\rm fixed}$ bytes, the corresponding memory models are
\[
M_{\rm single}=M_{\rm fixed}+M_{\rm block},
\qquad
\boxed{
M_{\rm double}=M_{\rm fixed}+2M_{\rm block}.
}
\]
The proposition is a scheduling model, not a hardware timing claim.  PCIe or
NVLink contention, kernel launch overhead, synchronization, page pinning, and
library workspaces can reduce realized overlap.

\begin{figure}[htbp]
\centering
\includegraphics[width=.82\linewidth]{figure_12_double_buffer_model.png}
\caption{Idealized effect of asynchronous double buffering over $64$ streamed blocks.  The largest modeled speedup occurs when transfer and computation times are comparable, because neither resource is left idle for most of the pipeline.  The additional cost is one extra resident block buffer.}
\label{fig:double-buffer-model}
\end{figure}


